%==============================================================================
%  APPENDIX - Detailed technology stack and justified comparison
%  Concrete implementations of the imposed foundation + evaluated alternatives
%==============================================================================

%==============================================================================
\chapter{Appendix: Technology stack and comparison of alternatives}
\label{chap:technologies}
%==============================================================================
The technical foundation of \textbf{Zümm} (MVC, Servlet, JSP, EJB, JDBC, XML
i18n, Google OAuth, X.509 / SSL, XML web service) is \textbf{imposed} by the
requirements specification and by the grading grid. The freedom of choice
therefore concerns the \textbf{concrete implementations} of this foundation:
application server, DBMS, front-end libraries, web-service style. This appendix
lists the technologies chosen layer by layer, then justifies each choice against
its alternatives.

%------------------------------------------------------------------------------
\section{Chosen technology stack}
%------------------------------------------------------------------------------
The table below specifies, for each layer of the multi-tier architecture
(figure~\ref{fig:couches}), the concrete implementation chosen and its reference
version. All the building blocks are \textbf{free and open source}, in line with
the requirement of self-deployment without a subscription.

\begin{xltabular}{\textwidth}{>{\bfseries}p{3.1cm} p{4.9cm} X}
\toprule
\textbf{Layer / need} & \textbf{Chosen technology} & \textbf{Role in Zümm} \\
\midrule
\endhead
Platform          & Jakarta EE 10 on JDK 17 LTS       & Current ``J2EE'' foundation: Servlet, JSP, EJB, JAX-WS \\
Application server & WildFly (Web + EJB container)     & Runs the WAR/EAR archive, provides the containers \\
Control (MVC)     & Jakarta Servlets + JSP + JSTL      & Request routing, control layer \\
Business          & Session EJB (stateless)           & Management rules: ROI, thresholds, honey quantity \\
Data access       & JDBC + HikariCP pool              & SQL queries, connection management \\
DBMS              & PostgreSQL (+ PostGIS)            & Relational persistence, site geolocation \\
Presentation      & HTML5 + JavaScript + Bootstrap 5  & Responsive UI, web / mobile parity \\
Charts            & Chart.js                          & Production, alerts and summary dashboards \\
Mapping           & Leaflet + OpenStreetMap           & Site map and foraging radii \\
i18n              & XML files + \texttt{ResourceBundle} & FR / EN / AR texts, Arabic RTL support \\
Configuration     & \texttt{ConfigZumm.ini} via Singleton  & Thresholds, units and modules without redeployment \\
Third-party web service & JAX-WS (SOAP / XML)         & \texttt{getZummHoneyActualQuantity(int)} \\
Authentication    & Google OAuth 2.0                  & Delegated login without password management \\
Exchange security & TLS / SSL + X.509 certificates    & Encryption and inter-component authentication \\
Offline           & PWA (Service Worker + IndexedDB)  & Field entry without a network, deferred synchronisation \\
Build \& tests    & Maven, JUnit 5, Mockito, Arquillian & Compilation, unit and in-container tests \\
\bottomrule
\end{xltabular}

\begin{encadre}
\textbf{Note:} moving from the historical name ``J2EE'' to \textbf{Jakarta EE 10}
changes neither the concepts nor the APIs noted (Servlet, JSP, EJB, JDBC). It
merely adopts the maintained version of the platform. The modernisation
trajectory towards Spring Boot is treated separately in
appendix~\ref{chap:migration}.
\end{encadre}

%------------------------------------------------------------------------------
\section{Justified comparison of the alternatives}
%------------------------------------------------------------------------------
For each open decision, the credible options were evaluated. The table keeps the
option, cites the discarded alternatives and gives the justification with regard
to the project requirements (self-deployment, free of charge, usability,
geolocation, interoperability).

\subsection{Application server}
\begin{xltabular}{\textwidth}{>{\bfseries}p{2.7cm} p{4.2cm} X}
\toprule
\textbf{Option} & \textbf{Discarded alternatives} & \textbf{Justification of the choice} \\
\midrule
\endhead
WildFly & GlassFish, Payara, Open~Liberty, Apache~TomEE & Complete and free \textbf{EJB + Servlet} container, fast startup, abundant documentation. TomEE remains lighter but less integrated. Payara also fits (production-oriented GlassFish fork). \\
\bottomrule
\end{xltabular}

\subsection{Database management system}
\begin{xltabular}{\textwidth}{>{\bfseries}p{2.7cm} p{4.2cm} X}
\toprule
\textbf{Option} & \textbf{Discarded alternatives} & \textbf{Justification of the choice} \\
\midrule
\endhead
PostgreSQL & MySQL / MariaDB, H2, Oracle XE & Native \textbf{geolocation} (PostGIS) for the sites and foraging radii, good support of \texttt{CHECK} constraints (frame/super rules), time series of measurements. Free and without proprietary lock-in, unlike Oracle. \\
\bottomrule
\end{xltabular}

\subsection{Third-party web-service style}
\begin{xltabular}{\textwidth}{>{\bfseries}p{2.7cm} p{4.2cm} X}
\toprule
\textbf{Option} & \textbf{Discarded alternatives} & \textbf{Justification of the choice} \\
\midrule
\endhead
JAX-WS (SOAP / XML) & JAX-RS (REST/JSON), gRPC & The requirements specification imposes an \textbf{XML web service} exposing \texttt{float getZummHoneyActualQuantity(int zummID)}. SOAP is its literal translation, with a WSDL contract queryable from Excel. A REST/JSON endpoint may complete the offer during modernisation (appendix~\ref{chap:migration}). \\
\bottomrule
\end{xltabular}

\subsection{Mapping and foraging radii}
\begin{xltabular}{\textwidth}{>{\bfseries}p{2.7cm} p{4.2cm} X}
\toprule
\textbf{Option} & \textbf{Discarded alternatives} & \textbf{Justification of the choice} \\
\midrule
\endhead
Leaflet + OSM & Google Maps JS API, Mapbox & \textbf{Free and without a paid key} or billed quota, consistent with the self-deployment requirement. Simple drawing of the foraging circles (1, 2, 3~km) from the site coordinates. \\
\bottomrule
\end{xltabular}

\subsection{Charts library}
\begin{xltabular}{\textwidth}{>{\bfseries}p{2.7cm} p{4.2cm} X}
\toprule
\textbf{Option} & \textbf{Discarded alternatives} & \textbf{Justification of the choice} \\
\midrule
\endhead
Chart.js & D3.js, ApexCharts, Highcharts & Low learning curve and direct rendering of the histograms and curves of the dashboards (production, ROI). D3 is more powerful but oversized. Highcharts is not free for commercial use. \\
\bottomrule
\end{xltabular}

\subsection{Offline and mobile access}
\begin{xltabular}{\textwidth}{>{\bfseries}p{2.7cm} p{4.2cm} X}
\toprule
\textbf{Option} & \textbf{Discarded alternatives} & \textbf{Justification of the choice} \\
\midrule
\endhead
PWA (Service Worker + IndexedDB) & Native application (Android/iOS), Flutter, React~Native & \textbf{Web / mobile parity} with a single codebase, field entry without a network then deferred synchronisation. Avoids the development and distribution of native applications, out of the academic scope. \\
\bottomrule
\end{xltabular}

\begin{encadre}
\textbf{Consistency with the academic constraint:} all the chosen building
blocks are free, documented and \emph{justifiable} ``why / how / where''
(\emph{Technologies} criterion of the grading grid, schedule chapter). None
relies on a code generator: they are implemented by hand by the team, in
compliance with the exam constraint.
\end{encadre}
